\section{The Risk That Does Not Design Out}
\label{sec:12.3}
Everything this book has argued for can be built correctly and still go wrong in ways nobody designed for. Chapter~\ref{sec:6} works through what that has actually looked like: the failure a red team finds before release, the failure an outside party finds because the developer did not look, the identified harm that then has to be remediated, and the adversary causing the harm on purpose. What those cases have in common is more useful at this distance than any of them is on its own.

In none of them did the fix come from the party that had promised to behave well. It came from somebody holding a specific lever: a research team with access to the vendor's own data, a set of universities exercising a purchasing decision, a company deciding a blunt technical limit was cheaper than moderating case by case, an organization that ran the test nobody inside had run. Set that beside what a decade of voluntary commitment produced. The instrument that separates the two is what a party has already paid, rather than what it has undertaken to pay.

That leaves one residual risk this book cannot dispose of, and it is the one chapter~\ref{sec:3} is about. Every lever above belongs to somebody, and a government that can reach the party holding the lever has the lever. Anticipation, remediation, and technical resilience each raise the price of misuse, and none of them changes who can be held in contempt for refusing an order. This is why the argument ended where it did. A constraint that holds against the party in possession has to be held by something that party cannot instruct, and the one version of that with a working instance behind it is a bearer with a stake of its own — still in the party's possession, which is why it has to be paired with custody the party does not control — at the price section~\ref{sec:3.8} sets out. Section~\ref{sec:3.3} puts two untried constructions ahead of it, and a deployment that goes straight to the bearer owes an account of what it tried first.

What the project is for, and how anyone would know it was working, are the questions this chapter has already answered. This is the part that stays open.
